\documentclass[12pt, a4paper]{report}

% Packages
\usepackage[utf8]{inputenc}
\usepackage[T1]{fontenc}
\usepackage[english]{babel}
\usepackage{amsmath, amssymb, amsthm, amsfonts}
\usepackage{geometry}
\usepackage{hyperref}
\usepackage{bookmark}
\usepackage{enumitem}
\usepackage{fancyhdr}
\usepackage{titlesec}
\usepackage{xcolor}
\usepackage{tikz}
\usepackage{graphicx}
\usepackage{mathrsfs} % For script fonts if needed

% Geometry setup
\geometry{
    top=2.5cm,
    bottom=2.5cm,
    left=2.5cm,
    right=2.5cm
}
\setlength{\headheight}{30pt}

% Hyperref setup
\hypersetup{
    colorlinks=true,
    linkcolor=blue,
    filecolor=magenta,      
    urlcolor=cyan,
    pdftitle={Stochastic Finance and Risk Modelling - Chapter 4},
    pdfauthor={Laurence Carassus and Gaoyue Guo},
}

% Theorem environments
\newtheorem{theorem}{Theorem}[chapter]
\newtheorem{proposition}[theorem]{Proposition}
\newtheorem{lemma}[theorem]{Lemma}
\newtheorem{corollary}[theorem]{Corollary}

\theoremstyle{definition}
\newtheorem{definition}{Definition}[chapter]
\newtheorem{example}{Example}[chapter]
\newtheorem{remark}[theorem]{Remark}

% Layout/Style
\pagestyle{fancy}
\fancyhf{}
\rhead{\leftmark}
\lhead{}
\cfoot{\thepage}

\title{\textbf{Stochastic Finance and Risk Modelling} \\ \large Chapter 4: Optimal portfolios}
\author{Based on slides by Romain Perchet \\ Laurence Carassus and Gaoyue Guo \\ \textit{Expanded Study Resource}}
\date{\today}

\begin{document}

\maketitle

\tableofcontents

\chapter*{Introduction}
\addcontentsline{toc}{chapter}{Introduction}
This document serves as a comprehensive study guide for the course ``Stochastic Finance and Risk Modelling''. It expands upon the lecture slides relative to Chapter 4: ``Optimal portfolios''.

In this chapter, we transition from the seller's perspective (which focuses on fair pricing through replication and risk-neutral measures) to the buyer's (investor's) perspective. The primary objective is to optimize a portfolio by maximizing the expected reward while controlling the associated risk. We define the reward and the variance of a portfolio and introduce the fundamental problems of optimal asset allocation. We first explore the Mean-Variance framework introduced by Harry Markowitz for a market containing only risky assets, leading to the concept of the Efficient Frontier. We then incorporate a riskless asset to define the general market portfolio and the Capital Market Line.

\chapter{Reward and Risk}

\section{From Seller to Buyer Perspective}

In the previous chapters, our analysis was framed within the context of the \textbf{seller} in a one-period model ($T=1$). A seller (e.g., a financial institution offering a derivative) aims to establish a ``fair'' pricing framework for trading financial products. This is achieved through the no-arbitrage condition, which guarantees the existence of a risk-neutral measure $\mathbb{Q}$ equivalent to the physical measure $\mathbb{P}$ ($\mathbb{Q} \sim \mathbb{P}$). Under $\mathbb{Q}$, all discounted asset prices are martingales, and the price of any contingent claim is given by its expected discounted payoff. The seller's primary goal is replication and hedging, seeking to eliminate risk entirely by replicating the payoff of the product sold.

In this chapter, we shift our focus to the \textbf{buyer} (or investor). The buyer's objective is fundamentally different: they aim to build a portfolio of assets to optimize their ``utility.'' This utility is generally driven by two conflicting desires:
\begin{enumerate}
    \item \textbf{Maximizing Reward}: Achieving the highest possible return on the investment.
    \item \textbf{Minimizing Risk}: Reducing the uncertainty or volatility associated with the return.
\end{enumerate}
Thus, optimal portfolio selection involves finding the best possible trade-off between reward and risk under the true probability measure $\mathbb{P}$, which models the actual observed behavior of the market.

\section{Portfolio Construction and Weights}

We consider a one-period financial market consisting of:
\begin{itemize}
    \item One riskless asset (e.g., a bank account or zero-coupon bond), denoted by $S^0 = (S_t^0)_{t=0,1}$. Its return is deterministic and denoted by $r$, such that $S_1^0 = S_0^0(1+r)$.
    \item $d$ risky assets (e.g., stocks), denoted by $S = (S^1, \dots, S^d) = (S_t^1, \dots, S_t^d)_{t=0,1}$. These assets are modeled as random variables on a probability space $(\Omega, \mathcal{F}, \mathbb{P})$.
\end{itemize}

A typical portfolio is constructed at time $t=0$ and held until $t=1$. It is entirely characterized by a trading strategy $(\xi^0, \xi) \in \mathbb{R} \times \mathbb{R}^d$, where:
\begin{itemize}
    \item $\xi^0$ is the number of units held in the riskless asset $S^0$.
    \item $\xi^i$ is the number of units held in the $i$-th risky asset $S^i$, for $i = 1, \dots, d$.
\end{itemize}

The \textbf{value of the portfolio} at time $t \in \{0, 1\}$ is simply the sum of the values of the individual holdings:
\begin{equation}
    V_t^{\xi^0, \xi} := \xi^0 S_t^0 + \xi \cdot S_t = \xi^0 S_t^0 + \sum_{i=1}^d \xi^i S_t^i, \quad t = 0, 1.
\end{equation}
Here, $V_0^{\xi^0, \xi}$ represents the initial wealth required to set up the portfolio.

It is often more convenient and intuitive to work with portfolio proportions, or \textbf{weights}, rather than the absolute number of shares ($\xi^i$). Assuming the initial wealth is non-zero ($V_0^{\xi^0, \xi} \neq 0$), we define the weight $w^i$ as the fraction of the initial wealth invested in asset $i$:
\begin{equation}
    w^i := \frac{\xi^i S_0^i}{V_0^{\xi^0, \xi}} = \frac{\xi^i S_0^i}{\sum_{j=0}^d \xi^j S_0^j} \quad \text{for } i = 0, \dots, d.
\end{equation}

By definition, the sum of all weights must equal $1$, representing $100\%$ of the initial wealth:
\begin{equation}
    \sum_{i=0}^d w^i = 1.
\end{equation}

\begin{remark}[Interpretation of Weights]
\leavevmode
\begin{itemize}
    \item If $w^i > 0$, the investor has a \textbf{long position} in asset $i$; they have purchased it.
    \item If $w^i < 0$, the investor has a \textbf{short position} in asset $i$; they have borrowed and sold it, hoping for its price to drop. The cash generated from this short sale can be used to invest in other assets.
\end{itemize}
\end{remark}

Using the weights, the portfolio value at time $t$ can be rewritten as:
\begin{equation}
    V_t^{\xi^0, \xi} = \sum_{i=0}^d \xi^i S_t^i = \sum_{i=0}^d \left( \frac{w^i V_0^{\xi^0, \xi}}{S_0^i} \right) S_t^i = V_0^{\xi^0, \xi} \sum_{i=0}^d w^i \frac{S_t^i}{S_0^i}, \quad t = 0, 1.
\end{equation}

Let $w := (w^1, \dots, w^d) \in \mathbb{R}^d$ be the vector of weights for the $d$ \textit{risky} assets. Let $\mathbf{1} := (1, \dots, 1)^T \in \mathbb{R}^d$ be the vector of ones. Then, the weight of the riskless asset can be expressed directly as a function of the risky portfolio weights:
\begin{equation}
    w^0 = 1 - \sum_{i=1}^d w^i = 1 - w \cdot \mathbf{1}.
\end{equation}
Because the entire portfolio composition is fully determined by the vector $w$ (since $w^0$ is dependent on it), we will henceforth denote the portfolio simply by $V^w$ instead of $V^{\xi^0, \xi}$.

\section{Portfolio Return Formulation}

The performance of a portfolio over the period $[0,1]$ is measured by its \textbf{return}. For a generic asset or portfolio $X$ with prices $X_0 \neq 0$ and $X_1$, its return is defined as:
\begin{equation}
    \rho(X) := \frac{X_1 - X_0}{X_0} = \frac{X_1}{X_0} - 1.
\end{equation}

Let us derive the return of the portfolio $V^w$. Using the definitions of the weights, we have:
\begin{align}
    \rho(V^w) &= \frac{V_1^w - V_0^w}{V_0^w} \notag \\
    &= \frac{\xi^0 (S_1^0 - S_0^0) + \sum_{i=1}^d \xi^i (S_1^i - S_0^i)}{V_0^w} \notag \\
    &= \frac{\xi^0 S_0^0}{V_0^w} \left( \frac{S_1^0 - S_0^0}{S_0^0} \right) + \sum_{i=1}^d \frac{\xi^i S_0^i}{V_0^w} \left( \frac{S_1^i - S_0^i}{S_0^i} \right) \notag \\
    &= w^0 \rho(S^0) + \sum_{i=1}^d w^i \rho(S^i).
\end{align}

Recall that the return on the riskless asset is defined as $r$. For consistency with the notation in the subsequent formulas, let's designate the return of the riskless asset as $R^0 := \rho(S^0) = r$.
Furthermore, let $\rho(S) := (\rho(S^1), \dots, \rho(S^d))^T$ be the random vector of returns for the $d$ risky assets. The portfolio return can thus be written compactly as a linear combination of individual asset returns:
\begin{equation}
    \rho(V^w) = w^0 R^0 + w \cdot \rho(S) = (1 - w \cdot \mathbf{1}) R^0 + w \cdot \rho(S).
\end{equation}
This demonstrates that the return of a portfolio is exactly the weighted sum of the returns of its constituent assets.

\section{Quantifying Reward and Risk}

To formulate an optimization problem, we need mathematical proxies for the investor's objectives (reward) and constraints/penalties (risk).

\subsection{Expected Return (Reward)}
The ``reward'' of a portfolio is universally measured by its \textbf{expected return} under the real-world probability measure $\mathbb{P}$. We denote the expected return of the $i$-th risky asset as $R_i$:
\begin{equation}
    R_i := \mathbb{E}^{\mathbb{P}}[\rho(S^i)] = \mathbb{E} \left[ \frac{S_1^i}{S_0^i} - 1 \right] \quad \text{for } i = 1, \dots, d.
\end{equation}
Let $R := (R_1, \dots, R_d)^T \in \mathbb{R}^d$ be the vector of expected returns for the risky assets. \textit{We assume that the vector $R$ is not collinear to the vector of ones $\mathbf{1}$ (i.e., not all risky assets have the exact same expected return, which avoids trivial uninteresting cases).}

The expected return of the total portfolio $V^w$, denoted $r(V^w)$, follows from the linearity of the expectation operator:
\begin{align}
    r(V^w) &:= \mathbb{E}[\rho(V^w)] \notag \\
    &= \mathbb{E}[(1 - w \cdot \mathbf{1}) R^0 + w \cdot \rho(S)] \notag \\
    &= (1 - w \cdot \mathbf{1}) R^0 + w \cdot \mathbb{E}[\rho(S)] \notag \\
    &= \mathbf{(1 - w \cdot \mathbf{1}) R^0 + w \cdot R}. \label{eq:expected_return}
\end{align}

\subsection{Variance (Risk)}
In the classical Markowitz framework, ``risk'' is quantified by the \textbf{variance} of the portfolio's return. The variance measures the degree of dispersion or volatility away from the expected return. A higher variance implies greater uncertainty in the final payoff.

Since the riskless asset return $R^0$ is a deterministic constant, its variance is zero, and it has zero covariance with all other assets. Therefore, the total variance of the portfolio depends exclusively on the risky assets' part: $\operatorname{Var}(\rho(V^w)) = \operatorname{Var}(w \cdot \rho(S) + \text{constant}) = \operatorname{Var}(w \cdot \rho(S))$.

Let $\Sigma$ be the $d \times d$ \textbf{covariance matrix} of the risky asset returns. Its elements $\sigma_{i,j}$ are given by:
\begin{equation}
    \sigma_{i,j} := \operatorname{Cov}(\rho(S^i), \rho(S^j)) = \frac{\operatorname{Cov}(S_1^i, S_1^j)}{S_0^i S_0^j}, \quad \text{for } 1 \leq i, j \leq d.
\end{equation}
Note that the diagonal elements $\sigma_{i,i}$ correspond to the variances of the individual assets' returns ($\operatorname{Var}(\rho(S^i))$).

The variance of the portfolio return can be calculated using the foundational properties of variance for linear combinations of random variables:
\begin{align}
    \operatorname{Var}(\rho(V^w)) &:= \mathbb{E}\left[ \left( \rho(V^w) - \mathbb{E}[\rho(V^w)] \right)^2 \right] \notag \\
    &= \operatorname{Var}\left( \sum_{i=1}^d w^i \rho(S^i) \right) \notag \\
    &= \sum_{i=1}^d \sum_{j=1}^d w^i w^j \operatorname{Cov}(\rho(S^i), \rho(S^j)) \notag \\
    &= \mathbf{w^T \Sigma w}. \label{eq:portfolio_variance}
\end{align}
This fundamental quadratic form, $w^T \Sigma w$, sits at the heart of portfolio optimization.

\vspace{0.3cm}
\textbf{Assumption on $\Sigma$:} We assume that the covariance matrix $\Sigma$ is invertible and \textit{positive definite}. This means that for any non-zero allocation vector $w \neq \mathbf{0}$, the portfolio variance is strictly positive:
\begin{equation*}
    w \neq \mathbf{0} \implies w^T \Sigma w > 0.
\end{equation*}
Economically, this implies that it is impossible to construct a completely riskless portfolio using only the risky assets, assuming a non-trivial position is taken. In real markets, true zero-risk combinations of purely risky assets do not exist.

\section{The General Optimization Objectives}

With mathematical descriptions for reward and risk, the investor's problem can be formally modeled. The problem involves a trade-off: higher returns generally require accepting higher risks. We define a subset $E \subseteq \mathbb{R}^d$ which represents the feasible region for trading weights (e.g., $E = \mathbb{R}^d$ for unconstrained trading including short selling, or $E = \mathbb{R}^d_+$ if short selling is forbidden).

There are three equivalent ways to formulate the mathematical optimization problem, yielding identical portfolios along the ``Efficient Frontier''.

\paragraph{Formulation 1: Maximize Reward subject to a Risk Constraint}
Let $\sigma > 0$ be the maximum acceptable portfolio variance (the risk budget). The investor seeks to maximize their expected return without exceeding this risk budget:
\begin{equation}
    \max_{w \in E} \left( (1 - w \cdot \mathbf{1}) R^0 + w \cdot R \right) \quad \text{subject to} \quad w^T \Sigma w \leq \sigma^2.
\end{equation}

\paragraph{Formulation 2: Minimize Risk subject to a Reward Constraint}
Let $\mu \in \mathbb{R}$ be the desired target portfolio return. The investor seeks the safest possible way (lowest variance) to achieve at least this expected return:
\begin{equation}
    \min_{w \in E} w^T \Sigma w \quad \text{subject to} \quad (1 - w \cdot \mathbf{1}) R^0 + w \cdot R \geq \mu.
\end{equation}

\paragraph{Formulation 3: Maximize a Risk-Adjusted Reward (The Utility Formulation)}
In modern finance, investors are often modeled using a combined objective function. Let $\lambda \ge 0$ denote the investor's \textbf{risk aversion parameter}. 
\begin{itemize}
    \item If $\lambda = 0$, the investor is risk-neutral and only cares about maximizing expected return.
    \item If $\lambda \gg 1$, the investor is highly risk-averse, punishing variance heavily.
\end{itemize}
The unconstrained optimization problem seeks to maximize a penalized expected return function:
\begin{equation}
    \max_{w \in E} \left( (1 - w \cdot \mathbf{1}) R^0 + w \cdot R - \frac{\lambda}{2} w^T \Sigma w \right).
\end{equation}
The term $\frac{\lambda}{2} w^T \Sigma w$ represents the utility penalty exacted by risk. Varying $\lambda$ from $0$ to $\infty$ maps out the entire set of optimal portfolios.

\chapter{The Market Portfolio (Purely Risky Assets)}

In this chapter, we restrict the investor's universe to only the $d$ risky assets. This means no cash is held in the bank account, or more formally, the weight assigned to the riskless asset is strictly zero:
\begin{equation}
    w^0 = 0.
\end{equation}
Because the sum of all portfolio weights must equal $1$, the condition $w^0 = 0$ is exactly equivalent to requiring that the weights of the risky assets sum up to $1$:
\begin{equation}
    w \cdot \mathbf{1} = \sum_{i=1}^d w^i = 1.
\end{equation}

Thus, the feasible set of trading strategies for the risky assets is constrained to the hyperplane $E_0$:
\begin{equation}
    E = E_0 := \{ w \in \mathbb{R}^d : w \cdot \mathbf{1} = 1 \}.
\end{equation}

Under this restriction, the three optimization problems formulated in the previous chapter specialize as follows:

\paragraph{Formulation 1 (Variance Target $\sigma^2$):}
\begin{equation}
    f(\sigma) := \max_{w \in E_0} w \cdot R \quad \text{subject to} \quad w^T \Sigma w \leq \sigma^2. \label{eq:p1_market}
\end{equation}

\paragraph{Formulation 2 (Return Target $\mu$):}
\begin{equation}
    g(\mu) := \min_{w \in E_0} w^T \Sigma w \quad \text{subject to} \quad w \cdot R \geq \mu. \label{eq:p2_market}
\end{equation}

\paragraph{Formulation 3 (Utility with Risk Aversion $\lambda$):}
\begin{equation}
    \max_{w \in E_0} \left( w \cdot R - \frac{\lambda}{2} w^T \Sigma w \right). \label{eq:p3_market}
\end{equation}

These formulations characterize the \textbf{Mean-Variance Efficient Frontier}. The efficient frontier represents the set of optimal portfolios that offer the highest possible expected return for a given level of risk, or conversely, the lowest possible risk for a given level of expected return. Portfolios lying strictly below the frontier are considered sub-optimal because superior alternatives exist.

\section{Optimization using Lagrange Multipliers}

We will explicitly solve for the optimal weights using Formulation 2 via the method of Lagrange multipliers. To build intuition, we first solve an unconstrained version (omitting the $w \cdot \mathbf{1}=1$ requirement), and then tackle the fully constrained problem.

\subsection{Unconstrained Optimization}
Suppose temporarily that we do not enforce $w \cdot \mathbf{1} = 1$ (which means $w^0$ is arbitrary, functioning as a slack variable, but without assigning any explicit return to it). We wish to minimize variance while achieving a target return $\mu$:
\begin{equation*}
    \min_{w \in \mathbb{R}^d} w^T \Sigma w \quad \text{subject to} \quad w \cdot R = \mu.
\end{equation*}
Note we use the equality constraint $w \cdot R = \mu$ instead of the inequality constraint, as achieving the absolute minimum risk will naturally bind the return to exactly the target $\mu$.

We introduce the \textbf{Lagrangian} with a single multiplier $\lambda$ corresponding to the target return constraint:
\begin{equation}
    L(w, \lambda) := w^T \Sigma w + \lambda(\mu - w \cdot R).
\end{equation}

By Fermat's theorem, optimal solutions must satisfy the First-Order Conditions (taking partial derivatives and equating to zero):
\begin{align}
    \nabla_w L(w, \lambda) &= 2\Sigma w - \lambda R = \mathbf{0}, \label{eq:foc_unconst_w} \\
    \frac{\partial L}{\partial \lambda} &= \mu - w \cdot R = 0. \label{eq:foc_unconst_lambda}
\end{align}
Equation \eqref{eq:foc_unconst_w} yields an expression for $w$ in terms of $\lambda$:
\begin{equation}
    2\Sigma w = \lambda R \implies w = \frac{\lambda}{2} \Sigma^{-1} R. \label{eq:w_unconstrained}
\end{equation}
Substitute \eqref{eq:w_unconstrained} into the return constraint \eqref{eq:foc_unconst_lambda}:
\begin{equation*}
    \mu = w \cdot R = w^T R = \left( \frac{\lambda}{2} \Sigma^{-1} R \right)^T R = \frac{\lambda}{2} R^T (\Sigma^{-1})^T R.
\end{equation*}
Since $\Sigma$ is symmetric, its inverse $\Sigma^{-1}$ is also symmetric, so $(\Sigma^{-1})^T = \Sigma^{-1}$:
\begin{equation}
    \mu = \frac{\lambda}{2} R^T \Sigma^{-1} R \implies \frac{\lambda}{2} = \frac{\mu}{R^T \Sigma^{-1} R}.
\end{equation}

Finally, substituting $\frac{\lambda}{2}$ back into \eqref{eq:w_unconstrained} gives the optimal portfolio weights $\tilde{w}$ for this unconstrained setting:
\begin{equation}
    \tilde{w} = \mu \frac{\Sigma^{-1} R}{R^T \Sigma^{-1} R}.
\end{equation}

\subsection{Constrained Optimization (E = E\_0)}
We now solve the actual problem defined in \eqref{eq:p2_market}, restricting the riskless asset weight to zero by enforcing $w \cdot \mathbf{1} = 1$:
\begin{equation}
    \min_{w \in \mathbb{R}^d} w^T \Sigma w \quad \text{subject to} \quad w \cdot R = \mu \quad \text{and} \quad w \cdot \mathbf{1} = 1.
\end{equation}

The \textbf{Lagrangian} now requires two multipliers: $\lambda$ for the expected return constraint and $\theta$ for the portfolio budget constraint:
\begin{equation}
    L(w, \lambda, \theta) := w^T \Sigma w + \lambda(\mu - w \cdot R) + \theta(1 - w \cdot \mathbf{1}).
\end{equation}

The First-Order Conditions are:
\begin{align}
    \nabla_w L &= 2\Sigma w - \lambda R - \theta \mathbf{1} = \mathbf{0}, \label{eq:foc_const_w} \\
    \frac{\partial L}{\partial \lambda} &= \mu - w \cdot R = 0, \label{eq:foc_const_lambda} \\
    \frac{\partial L}{\partial \theta} &= 1 - w \cdot \mathbf{1} = 0. \label{eq:foc_const_theta}
\end{align}

From \eqref{eq:foc_const_w}, we solve for $w$:
\begin{equation}
    2\Sigma w = \lambda R + \theta \mathbf{1} \implies \tilde{w} = \frac{1}{2}\lambda \Sigma^{-1} R + \frac{1}{2}\theta \Sigma^{-1} \mathbf{1}. \label{eq:w_constrained_param}
\end{equation}

To find the optimal multipliers $\tilde{\lambda}$ and $\tilde{\theta}$, we substitute $\tilde{w}$ into the two constraints \eqref{eq:foc_const_lambda} and \eqref{eq:foc_const_theta}. Before doing so, it is extremely useful to define three scalar constants that encapsulate matrix-vector products:
\begin{equation}
    a := R^T \Sigma^{-1} R, \quad b := R^T \Sigma^{-1} \mathbf{1} = \mathbf{1}^T \Sigma^{-1} R, \quad c := \mathbf{1}^T \Sigma^{-1} \mathbf{1}.
\end{equation}
These scalars $a, b,$ and $c$ solely depend on the market parameters (expected returns $R$ and covariance matrix $\Sigma$) and are thus fixed constants for the optimization problem.

Now substitute \eqref{eq:w_constrained_param} into constraint 1 ($R^T \tilde{w} = \mu$):
\begin{align*}
    R^T \left( \frac{1}{2}\lambda \Sigma^{-1} R + \frac{1}{2}\theta \Sigma^{-1} \mathbf{1} \right) &= \mu \\
    \frac{1}{2}\lambda (R^T \Sigma^{-1} R) + \frac{1}{2}\theta (R^T \Sigma^{-1} \mathbf{1}) &= \mu \\
    \frac{1}{2}\lambda a + \frac{1}{2}\theta b &= \mu. \quad \text{(Eq. A)}
\end{align*}

Substitute \eqref{eq:w_constrained_param} into constraint 2 ($\mathbf{1}^T \tilde{w} = 1$):
\begin{align*}
    \mathbf{1}^T \left( \frac{1}{2}\lambda \Sigma^{-1} R + \frac{1}{2}\theta \Sigma^{-1} \mathbf{1} \right) &= 1 \\
    \frac{1}{2}\lambda (\mathbf{1}^T \Sigma^{-1} R) + \frac{1}{2}\theta (\mathbf{1}^T \Sigma^{-1} \mathbf{1}) &= 1 \\
    \frac{1}{2}\lambda b + \frac{1}{2}\theta c &= 1. \quad \text{(Eq. B)}
\end{align*}

This gives a 2x2 linear system in terms of the variables $(\frac{\lambda}{2}, \frac{\theta}{2})$:
\begin{equation*}
    \begin{pmatrix} a & b \\ b & c \end{pmatrix} \begin{pmatrix} \lambda/2 \\ \theta/2 \end{pmatrix} = \begin{pmatrix} \mu \\ 1 \end{pmatrix}
\end{equation*}

Solving via Cramer's rule or direct substitution gives the optimal multipliers:
\begin{equation}
    \tilde{\lambda} = 2 \frac{\mu c - b}{ac - b^2}, \quad \quad \tilde{\theta} = 2 \frac{a - \mu b}{ac - b^2}. \label{eq:optimal_multipliers}
\end{equation}
The denominator is the determinant of the $2 \times 2$ matrix: $\Delta = ac - b^2$. By the Cauchy-Schwarz inequality applied to the positive definite metric given by $\Sigma^{-1}$, one can show that $ac > b^2$, so the determinant is always strictly positive ($ac - b^2 > 0$).

Substituting the expressions for $\tilde{\lambda}$ and $\tilde{\theta}$ \eqref{eq:optimal_multipliers} back into \eqref{eq:w_constrained_param} provides the exact optimal portfolio weights:
\begin{equation}
    \tilde{w} = \frac{\mu c - b}{ac - b^2} \Sigma^{-1} R + \frac{a - \mu b}{ac - b^2} \Sigma^{-1} \mathbf{1}. \label{eq:w_optimal_unsimplified}
\end{equation}

\section{The Two-Fund Separation Theorem}

We analyze the structure of the optimal portfolio $\tilde{w}$ derived above by grouping the terms that depend on the target return $\mu$ and those that do not.
Expanding \eqref{eq:w_optimal_unsimplified}:
\begin{align}
    \tilde{w} &= \frac{\mu c \Sigma^{-1} R - b \Sigma^{-1} R + a \Sigma^{-1} \mathbf{1} - \mu b \Sigma^{-1} \mathbf{1}}{ac - b^2} \notag \\
    &= \mu \left( \frac{c \Sigma^{-1} R - b \Sigma^{-1} \mathbf{1}}{ac - b^2} \right) + \left( \frac{a \Sigma^{-1} \mathbf{1} - b \Sigma^{-1} R}{ac - b^2} \right).
\end{align}

Let us define two specific, constant vectors $g$ and $h$ in $\mathbb{R}^d$:
\begin{equation}
    h := \frac{c \Sigma^{-1} R - b \Sigma^{-1} \mathbf{1}}{ac - b^2} \quad \text{and} \quad g := \frac{a \Sigma^{-1} \mathbf{1} - b \Sigma^{-1} R}{ac - b^2}.
\end{equation}
Crucially, vectors $h$ and $g$ consist solely of the market parameters $R$ and $\Sigma$; they are entirely independent of the investor's specific target return $\mu$.

Using these definitions, the optimal portfolio can be elegantly rewritten as an affine function of $\mu$:
\begin{equation}
    \tilde{w}(\mu) = \mu h + g. \label{eq:two_fund_separation}
\end{equation}

\subsection{Interpretation of g and h}
Equation \eqref{eq:two_fund_separation} is a powerful result known as the \textbf{Two-Fund Separation Theorem} (or Mutual Fund Theorem).
It states that \textit{every} optimal portfolio on the efficient frontier is merely a linear combination of two fixed portfolios: $g$ and $h$. Thus, an investor does not need to separately analyze and invest in all $d$ individual underlying stocks. They only need access to two ``mutual funds'' constructed according to the weights $g$ and $h$, respectively. They achieve their desired risk-return profile simply by scaling their investment amounts in these two funds based on their personal target $\mu$.

Specifically:
\begin{itemize}
    \item $g$ is exactly the portfolio weights required to achieve an expected return of $\mu = 0$ (set $\mu=0$ in the equation).
    \item $g+h$ represents the portfolio weights corresponding to an expected return of $\mu = 1$.
\end{itemize}

\section{Visualizing the Frontier: The Markowitz Bullet}

When we compute the minimum variance required $g(\mu)$ for various levels of $\mu$ and plot standard deviation $\sigma = \sqrt{g(\mu)}$ on the horizontal x-axis and expected return $\mu$ on the vertical y-axis, the resulting locus of points traces out a hyperbola in the variance-return ($\sigma^2$-$\mu$) space, or a hyperbolic shape in the standard deviation-return ($\sigma$-$\mu$) space known as the \textbf{Markowitz bullet}.

The left-most point on this curve denotes the \textbf{Global Minimum Variance Portfolio (GMVP)}. This portfolio achieves the absolute lowest possible risk across the entire trading universe $E_0$, without any specific constraint on return.

The frontier is split into two halves:
\begin{itemize}
    \item \textbf{The Efficient Frontier}: The upper limb of the hyperbola, above the GMVP. These portfolios offer the maximum return for a given level of risk. This is the only portion relevant to rational investors.
    \item \textbf{The Inefficient Frontier}: The lower limb below the GMVP. A rational investor would never choose a portfolio on this limb because there is always a corresponding vertical point on the upper limb offering higher returns for the exact same risk.
\end{itemize}

A sub-optimal generic portfolio (such as a totally un-diversified portfolio consisting of just a single stock, e.g. $w = (1, 0, 0)$) lies strictly to the right of the efficient frontier, inside the bullet envelope. This means the investor can improve their position by moving straight \textit{left} (minimizing variance while maintaining the same expected return via problem \eqref{eq:p2_market}) or moving straight \textit{up} (maximizing expected return while capping the variance at its current level via problem \eqref{eq:p1_market}) to reach the efficient frontier curve.

\chapter{The General Portfolio (Adding a Riskless Asset)}

We now return to the full unconstrained problem, allowing the investor's trading universe to include the \textbf{riskless asset} (such as a bank account or government bond) with a deterministic, guaranteed return rate $R^0 = r$. The addition of this asset fundamentally alters the shape of the efficient frontier and dramatically simplifies optimal portfolio selection.

Recall the full, multi-asset portfolio weights are defined by $w^0$ for the riskless asset and $w = (w^1, \dots, w^d) \in \mathbb{R}^d$ for the risky assets, such that $w^0 + w \cdot \mathbf{1} = 1$. The investor is no longer restricted to the hyperplane $E_0$.

Under this expanded universe, the \textbf{Formulation 2} optimization problem (minimizing variance for a target expected return $\mu$) becomes:
\begin{equation}
    \min_{(w^0, w) \in \mathbb{R} \times E} w^T \Sigma w \quad \text{subject to} \quad w^0 R^0 + w \cdot R \ge \mu, \quad \text{and} \quad w^0 + w \cdot \mathbf{1} = 1.
\end{equation}

Note the following crucial observations regarding the generalized problem:
\begin{itemize}
    \item The risk of the portfolio (the variance $w^T \Sigma w$) is still \textit{entirely determined} by the allocation vector $w$ of the risky assets. The riskless asset contributes zero variance, by definition.
    \item Any optimal portfolio will always be a simple two-asset linear combination of:
    \begin{enumerate}
        \item A single, optimally constructed portfolio composed entirely of risky assets.
        \item The riskless asset.
    \end{enumerate}
    \item Changing the target expected return $\mu$ merely alters the \textit{relative proportions} allocated between these two components, rather than changing the internal composition of the risky portfolio itself. 
\end{itemize}

\section{The Tangency Portfolio and the Market Portfolio}

To formalize this two-component separation property, we decompose the total portfolio into risky and riskless parts. Assume $w^0 < 1$ (the investor is not 100\% invested in the riskless asset, so they must hold some risky assets). Let us denote the interval of meaningful riskless weights as $I := (-\infty, 1)$.

For any allocation $w^0 \in I$, we can re-scale the risky asset weights to sum up to $1$. We define the normalized risky portfolio $\hat{w}$:
\begin{equation*}
    \hat{w} := \frac{w}{1 - w^0}.
\end{equation*}
By construction, this normalized portfolio satisfies $\hat{w} \cdot \mathbf{1} = \frac{w \cdot \mathbf{1}}{1 - w^0} = \frac{1 - w^0}{1 - w^0} = 1$. Therefore, $\hat{w}$ represents a purely risky portfolio that lies on the original (risky-only) efficient frontier defined in Chapter 2, i.e., $\hat{w} \in E_0$.

The expected return constraint can also be rewritten in terms of this normalized portfolio. We need $w^0 R^0 + w \cdot R = \mu$. Substituting $w = (1 - w^0)\hat{w}$ gives:
\begin{align*}
    w^0 R^0 + (1 - w^0)\hat{w} \cdot R &= \mu \\
    (1 - w^0)\hat{w} \cdot R &= \mu - w^0 R^0 \\
    \hat{w} \cdot R &= \frac{\mu - w^0 R^0}{1 - w^0}.
\end{align*}
Let us define the target return \textit{required from the risky portion} as:
\begin{equation}
    \hat{\mu}(w^0) := \frac{\mu - w^0 R^0}{1 - w^0}.
\end{equation}

The optimization problem can now be viewed sequentially: for any chosen level of cash $w^0$, we must find the optimal purely risky portfolio $\hat{w}$ that minimizes variance while yielding the return $\hat{\mu}(w^0)$.
Using the normalized weights, the variance objective function is:
\begin{equation}
    w^T \Sigma w = ((1-w^0)\hat{w})^T \Sigma ((1-w^0)\hat{w}) = (1-w^0)^2 \hat{w}^T \Sigma \hat{w}. \label{eq:var_normalized}
\end{equation}
Since $(1-w^0)^2$ is just a positive scalar constant for a fixed $w^0$, minimizing $(1-w^0)^2 \hat{w}^T \Sigma \hat{w}$ subject to $\hat{w} \cdot R \ge \hat{\mu}(w^0)$ is mathematically identical to solving the purely-risky Markowitz problem \eqref{eq:p2_market} for a target return of $\hat{\mu}(w^0)$!

Thus, the global optimization problem simplifies to:
\begin{align}
    \min_{w^0 \in I} \left[ \min_{\hat{w} \in E_0 : \hat{w} \cdot R \ge \hat{\mu}(w^0)} (1-w^0)^2 \hat{w}^T \Sigma \hat{w} \right]
    &= \min_{w^0 \in I} (1-w^0)^2 \left( \min_{\hat{w} \in E_0 : \hat{w} \cdot R \ge \hat{\mu}(w^0)} \hat{w}^T \Sigma \hat{w} \right) \notag \\
    &= \min_{w^0 \in I} \left[ (1-w^0)^2 g(\hat{\mu}(w^0)) \right],
\end{align}
where $g(\cdot)$ is exactly the variance-minimizing function for purely risky assets derived in Chapter 2. 
More precisely, for a given target expected return $m$, $g(m)$ represents the minimum variance achievable by a portfolio consisting \textit{only} of risky assets:
\begin{equation}
    g(m) := \min_{\hat{w} \in E_0 : \hat{w} \cdot R \ge m} \hat{w}^T \Sigma \hat{w}.
\end{equation}
This function $g$ describes the boundary of the purely risky efficient frontier (the Markowitz bullet).

This result implies that all \textit{globally} efficient portfolios (involving both riskless and risky assets) can be represented functionally as:
\begin{equation}
    \rho(V) = \tilde{w}^0 R^0 + (1 - \tilde{w}^0) \tilde{w} \cdot \rho(S).
\end{equation}
Here, $\tilde{w}$ denotes the fixed, unique optimal vector of risky asset weights (normalized so they sum to 1).

\textbf{Relationship to the purely risky case:} In Chapter 2 (without a risk-free asset), the optimal risky portfolio $w(\mu)$ changed depending on the investor's target return $\mu$, following the Two-Fund Separation Theorem ($w(\mu) = \mu h + g$). To achieve higher returns, the investor had to physically alter the composition of their risky assets (moving along the Markowitz bullet). 

However, introducing a riskless asset fundamentally changes this dynamic. The optimal relative composition of risky assets $\tilde{w}$ becomes \textit{completely independent} of the target return $\mu$. Regardless of whether an investor is highly conservative or aggressively seeking high returns, they should always hold exactly the same relative proportions of risky assets (the vector $\tilde{w}$). To achieve their specific target return $\mu$, they simply adjust $\tilde{w}^0$, which dictates the mix between cash and this single optimal risky portfolio (borrowing or lending at the risk-free rate).

In the literature, this uniquely optimal purely risky portfolio $\tilde{w}$ is called the \textbf{Market Portfolio} or the \textbf{Tangency Portfolio}, denoted $V_M$.
\begin{equation*}
    V_M := \tilde{w} \cdot S.
\end{equation*}

\section{The Capital Market Line (CML) and Conclusion}

Graphically mapping standard deviation (risk, $\sigma$) on the x-axis against expected return ($\mu$) on the y-axis reveals the profound impact of the riskless asset.

Because the riskless asset has zero variance and return $r = R^0$, it sits on the y-axis at $(0, r)$. An investor constructing a portfolio composed of fraction $w^0$ in the riskless asset and $(1-w^0)$ in any arbitrary risky portfolio $V_P$ will achieve an expected return and standard deviation that lie on a perfectly straight line connecting $(0, r)$ to the point representing $V_P$ structurally.

To maximize return for any given risk level, the investor must select the specific risky portfolio $V_P$ that creates the steepest possible slope originating from the riskless point $(0, r)$. Geometrically, the line with the maximum possible slope must be \textit{tangent} to the upper envelope of the purely risky efficient frontier (the Markowitz bullet). The single touchpoint where this straight line is tangent to the hyperbola represents the \textbf{Tangency Portfolio} or \textbf{Market Portfolio} $V_M$.

This tangent line dominates all other portfolios; we refer to this boundary as the \textbf{Capital Market Line (CML)}. The new Efficient Frontier is simply the CML itself. Every optimal strategy involves moving along this line by either:
\begin{itemize}
    \item \text{Lending (Holding cash):} Selecting a point on the line segment strictly between $(0, r)$ and $V_M$. The investor puts some money in the bank ($w^0 > 0$) and the rest into the exact combination of the market portfolio.
    \item \text{Borrowing (Using leverage):} Selecting a point on the line extending past $V_M$ to the upper-right. The investor borrows cash at the riskless rate ($w^0 < 0$) and invests more than $100\%$ of their initial capital into the market portfolio to amplify their returns.
\end{itemize}

\subsection{Equation of the Capital Market Line}

Let us derive the explicit formula for the CML. With the market portfolio defined as $V_M = \tilde{w} \cdot S$, any portfolio $V$ lying on the CML is a linear combination of the riskless asset and $V_M$:
\begin{equation}
    \rho(V) = \tilde{w}^0 \rho(S^0) + (1 - \tilde{w}^0) \rho(V_M).
\end{equation}
We evaluate the expected value (return) and the variance (risk) of this generic CML portfolio. First, the expected return $r(V)$:
\begin{equation}
    r(V) = \tilde{w}^0 R^0 + (1 - \tilde{w}^0) r(V_M) = R^0 + (1 - \tilde{w}^0) \big( r(V_M) - R^0 \big). \label{eq:cml_return}
\end{equation}
Second, the variance $\operatorname{Var}(\rho(V))$. Because the riskless asset has zero variance and zero covariance with $V_M$:
\begin{equation}
    \operatorname{Var}(\rho(V)) = (1 - \tilde{w}^0)^2 \operatorname{Var}(\rho(V_M)). \label{eq:cml_variance}
\end{equation}

By taking the square root of both sides, we can isolate the allocation term $(1 - \tilde{w}^0)$ as the ratio of standard deviations:
\begin{equation*}
    (1 - \tilde{w}^0) = \sqrt{ \frac{\operatorname{Var}(\rho(V))}{\operatorname{Var}(\rho(V_M))} } = \frac{\sigma_V}{\sigma_{M}}.
\end{equation*}

Substituting this ratio back into the expected return equation \eqref{eq:cml_return} yields the fundamental \textbf{Equation of the Capital Market Line}:
\begin{equation}
    r(V) = R^0 + \frac{r(V_M) - R^0}{\sqrt{\operatorname{Var}(\rho(V_M))}} \sqrt{\operatorname{Var}(\rho(V))}
\end{equation}
or, more simply in terms of standard deviations $\sigma$:
\begin{equation}
    \mu = R^0 + \frac{\mu_M - R^0}{\sigma_M} \cdot \sigma
\end{equation}

The slope of the line, $\frac{\mu_M - R^0}{\sigma_M}$, is known as the \textbf{Sharpe Ratio} of the market portfolio. It represents the investor's reward-to-risk ratio—the extra expected return earned per unit of standard deviation risk taken by moving away from the riskless baseline. Because the CML has the steepest possible slope, the tangency portfolio $V_M$ is precisely the specific blend of risky assets that provides the maximum possible Sharpe ratio.


\end{document}
